\clearpage
\VerdierPhysicalPageAnchor{0167}%
\begingroup\small\setlength{\parskip}{0pt}\setlength{\emergencystretch}{2em}
\VerdierSemanticAnchor{ch4-4-2-3-kernel-image}
On vérifie que :
\[
\operatorname{Ker}\bigl(\Psi^n(\omega_3)\bigr)
=Z\!\begin{bmatrix}n\\ f_3f_2,f_4\end{bmatrix},
\qquad
\operatorname{Im}\bigl(\Psi^n(\omega_3)\bigr)
=B\!\begin{bmatrix}n+1\\ f_2,f_4f_3\end{bmatrix}
 \Big/
 B\!\begin{bmatrix}n+1\\ f_2,f_3\end{bmatrix}.
\]
Le morphisme $\Psi^n(\omega_3)$ définit par suite un isomorphisme :
\VerdierSemanticAnchor{ch4-4-2-3-phi-isomorphism}
\[
\Phi^n(\omega_3)\colon
Z\!\begin{bmatrix}n\\ f_3,f_4\end{bmatrix}
 \Big/
Z\!\begin{bmatrix}n\\ f_3f_2,f_4\end{bmatrix}
\xrightarrow{\ \sim\ }
B\!\begin{bmatrix}n+1\\ f_2,f_4f_3\end{bmatrix}
 \Big/
B\!\begin{bmatrix}n+1\\ f_2,f_3\end{bmatrix}.
\]

Soit alors
$\omega_5=(\mathord\bullet\xrightarrow{f_1}\mathord\bullet
\xrightarrow{f_2}\mathord\bullet\xrightarrow{f_3}\mathord\bullet
\xrightarrow{f_4}\mathord\bullet\xrightarrow{f_5}\mathord\bullet)$
un objet de $\mathcal{F}\ell_5(\mathcal J)$. Les morphismes canoniques :
\[
E\!\begin{bmatrix}n\\ f_3,f_4,f_5\end{bmatrix}
\longrightarrow
Z\!\begin{bmatrix}n\\ f_3,f_4\end{bmatrix}
 \Big/
Z\!\begin{bmatrix}n\\ f_3f_2,f_4\end{bmatrix}
\longrightarrow0,
\]
\[
0\longrightarrow
B\!\begin{bmatrix}n+1\\ f_2,f_4f_3\end{bmatrix}
 \Big/
B\!\begin{bmatrix}n+1\\ f_2,f_3\end{bmatrix}
\longrightarrow
E\!\begin{bmatrix}n+1\\ f_1,f_2,f_3\end{bmatrix},
\]
composés avec l'isomorphisme $\Phi^n(f_2,f_3,f_4)$, définissent un
morphisme :
\VerdierSemanticAnchor{ch4-4-2-3-2}
\[
d^n(\omega_5)=d^n(f_1,f_2,f_3,f_4,f_5)\colon
E\!\begin{bmatrix}n\\ f_3,f_4,f_5\end{bmatrix}
\longrightarrow
E\!\begin{bmatrix}n+1\\ f_1,f_2,f_3\end{bmatrix}.
\tag{4.2.3.2}
\]
On a :
\VerdierSemanticAnchor{ch4-4-2-3-3}
\[
\left\{
\begin{aligned}
\operatorname{Im}\bigl(d^n(\omega_5)\bigr)
 &=B\!\begin{bmatrix}n+1\\ f_2,f_4f_3\end{bmatrix}
   \Big/
   B\!\begin{bmatrix}n+1\\ f_2,f_3\end{bmatrix},\\[2pt]
\operatorname{Ker}\bigl(d^n(\omega_5)\bigr)
 &=Z\!\begin{bmatrix}n\\ f_3f_2,f_4\end{bmatrix}
   \Big/
   B\!\begin{bmatrix}n\\ f_4,f_5\end{bmatrix}.
\end{aligned}
\right.
\tag{4.2.3.3}
\]
Le morphisme $d^n(\omega_5)$ dépend fonctoriellement, en un sens évident,
de $\omega_5$.

\VerdierSemanticAnchor{ch4-4-2-4}
\medskip\noindent\textbf{Proposition 4.2.4.}\quad
\textit{Soit
$\omega_7=(\mathord\bullet\xrightarrow{f_1}\mathord\bullet
\xrightarrow{f_2}\mathord\bullet\xrightarrow{f_3}\mathord\bullet
\xrightarrow{f_4}\mathord\bullet\xrightarrow{f_5}\mathord\bullet
\xrightarrow{f_6}\mathord\bullet\xrightarrow{f_7}\mathord\bullet)$
un objet de $\mathcal{F}\ell_7(\mathcal J)$. Posons :}
\[
d_1=d^{n-1}(f_3,f_4,f_5,f_6,f_7),
\qquad
d_2=d^n(f_1,f_2,f_3,f_4,f_5).
\]
\endgroup
